\emph{By Adam Turl}

\begin{flushleft}\emph{The German Communists are Communists because through all the intermediate stations and compromises, created not by them, but by the course of historical development, they clearly perceive and constantly pursue the final aim, viz., the abolition of classes and the creation of a society in which there will be no private ownership of land or of the means of production. The thirty-three Blanquists are Communists because they imagine that merely because they want to skip the intermediate stations and compromises, that settles the matter\ldots\ What childish innocence it is to present impatience as a theoretically convincing argument.}\\
--Frederick Engels, 1874.\footnote{V.\ I.\ Lenin, \emph{Left-Wing Communism: An Infantile Disorder} (New York: International Publishers, 1989), 37.}
\end{flushleft}

The Russian Revolution of October 1917 sent a shock through the international working class. Russia was only the first of a wave of revolutionary upheavals that gripped all of Europe in the wake of the First World War. British Prime Minister Lloyd George wrote that, ``the whole existing order in its political, social and economic aspects is questioned by the masses of the population from one end of Europe to the other.''\footnote{Quoted in John Rees, ``In Defense of October,'' \emph{International Socialism} 52, Autumn 1991, 9.}

Hundreds of thousands of radicalizing workers in Europe looked to the Bolshevik Revolution for inspiration and lessons. But in this first wave, only in Russia were workers able to take power. Only in Russia had revolutionaries organized themselves beforehand into a party independent of those socialists who favored compromise with the system and who opposed workers' power as ``premature.'' Only in Russia was there a party of revolutionary workers that had organized, agitated, and struggled unencumbered by the hesitancy of the reformists. Only in Russia had workers been able to overthrow the old bourgeois form of government---the duma, or parliament---and replace it with direct workers' democracy in the form of soviets (councils of elected workers' and soldiers' delegates).

That lesson spread quickly throughout the revolutionary workers' movement in Europe, keen to break away from the mass reformist social democratic parties like the Social Democratic Party (SPD) in Germany that had betrayed Marxist internationalism by siding with their own government's war effort. But revolutionaries around the world often failed to ask how the Bolsheviks had built that kind of party---a revolutionary party---over the years when most workers in Russia weren't revolutionaries yet. As Lenin put it, ``Only the history of Bolshevism during the whole period of its existence can satisfactorily explain why it was able to build up and to maintain under the most difficult conditions the iron discipline that is needed for the victory of the proletariat.''\footnote{Lenin, 10.}

Even before the revolutionary wave that ended the First World War began to ebb, the newly formed communist parties---organized along with the Russian Bolsheviks into the new Communist International, or Comintern---had a problem. While the number of organized revolutionary socialists was quite large---compared to today---they were still a minority of the workers' movement. The German Communist Party (KPD) for example, formed in December of 1918, had a membership barely over one hundred thousand members by the late summer of 1919, and most of them were not factory workers, but impatient radical youth, half of whom soon split to form a more ``left'' party that refused to work in existing trade unions or participate in bourgeois elections. In contrast, the reformist SPD had millions of members, and the ``centrist'' Independent Social Democrats (USP) had more than six hundred thousand.

The young revolutionaries were thrilled by the Russian Revolution, but not clear about the strategy and tactics employed to achieve it. For them, the organizational and tactical maneuvers required to bridge the gap between themselves and the majority of workers who had not yet made revolutionary conclusions was a non-question. For them, the simple equation---break with the treacherous reformists and prepare for revolution---was sufficient. It was for this reason that Lenin wrote \emph{Left-Wing Communism: An Infantile Disorder} for the Second Congress of the Comintern in 1920.

\sectiondivider

\paragraph{``Obsolete for whom?''}

In 1918 Germany was in the throes of revolution not unlike what had gripped Russia in early 1917. Soldiers mutinied and workers demonstrated and seized factories, and the whole apparatus of state power trembled. A national congress of workers and soldiers' councils formed, much like the Russian soviets. At the urging of the Social Democrats, the congress voted 344 to 98 to allow the election of a new national assembly. This was not an attempt to expand democratic rights after the overthrow of the German Kaiser, but, as in Russia, an attempt by the reformist SPD to save capitalism by urging the dissolution of direct workers' democracy in the councils. Workers had created organs of self-rule, and then voted to create a bourgeois assembly that would act to dismantle them. How to square this circle? The KPD at its founding conference voted sixty-two to twenty-three to boycott the elections.\footnote{Duncan Hallas, \emph{The Comintern}, (London, Chicago, Melbourne: Bookmarks, 1989), 38.} They argued, ``One must emphatically reject\ldots all reversion to parliamentary forms of struggle, which have become historically and politically obsolete.''\footnote{Lenin, 40.}

But as the British Marxist Duncan Hallas asked, ``Politically obsolete for whom?''\footnote{Hallas, 39.} For the millions of workers who were planning to vote? Lenin argued:

\blockquote{Parliamentarism has become ``historically obsolete.'' That is true as regards propaganda. But everyone knows that this is still a long way from overcoming it practically. Capitalism could have been declared, and quite rightly, to be ``historically obsolete'' many decades ago, but that doesn't at all remove the need for a very long and persistent struggle on the soil of capitalism.\footnote{Lenin, 40.}}

Being ``on the soil of capitalism'' didn't mean voting for anything that came along to the left of the Kaiser. It did, however, mean using the elections, which workers were looking to for change, to win workers from the SPD and to the idea of workers' control. As it turned out, more than eleven million German workers voted for the SPD that year.\footnote{Hallas, 39.}

The German Revolution, much like the February Revolution in Russia, ended the war, got rid of the monarchy, brought about workers' councils, and created a situation of ``dual power'' between a bourgeois government and a potential new workers' government in the workers' councils. But workers voted for the SPD, which had led them into war in the first place, betraying the socialist principles it claimed that it stood for. And the SPD was now trying to restore ``order'' and stop the revolution in its tracks. But the old tradition of looking to the SPD died hard and the KPD had yet to prove itself. The German ultra-lefts complained many workers were not only voting for reformists, but out and out reactionaries:

\blockquote{The millions of workers who still follow the Policy of the Centre [the Catholic ``Center'' Party] are counter-revolutionary. The rural proletarians provide legions of counter-revolutionary troops.\footnote{Lenin, 41.}}

But how could parliamentarism be obsolete Lenin asked, if

\blockquote{``Millions'' and ``legions'' of proletarians are not only still in favor of parliamentarism in general, but are downright ``counter-revolutionary''!? Clearly, parliamentarism in Germany is not yet politically obsolete. Clearly, the ``Lefts'' have mistaken their desire, their ideological-political attitude for actual fact. This is the most dangerous mistake revolutionaries can make.\footnote{Ibid.}}

The danger of this mistake had been most clearly revealed in the 1919 ``Spartakus Rising.'' The newly-formed Spartakusbund, only recently split away from the reformist SPD and comprised of a few thousand young and inexperienced militants, swung from their policy of abstaining from the elections to attempting to seize power on their own, before the majority of workers, even in Berlin, had been won over. The most experienced leaders of the party---Rosa Luxemburg and Karl Liebknecht---initially opposed the uprising but were overruled. The isolated insurrection was easily routed and the German Revolution lost two of its most capable leaders in the aftermath when Luxemburg and Liebknecht were murdered.\footnote{Chris Harman, \emph{The Lost Revolution} (Chicago: Haymarket Books, 2003), see pages 72--93.}

Since workers and peasants in Western Europe were ``more imbued with bourgeois-democratic and parliamentary prejudices than they were in Russia,'' Lenin argued that it was ``only from within such institutions as bourgeois parliaments that Communists can (and must) wage a long and persistent struggle to\ldots expose, dissipate and overcome these prejudices.''\footnote{Lenin, 48.} To really prepare for an insurrection and to make a revolution of the working class, it required using every means available to win over the workers who were wavering, still on the fence, or in the camp of the SPD. In Russia, under the tsar, the Bolsheviks had made regular use of elections in the Duma. As Tony Cliff wrote of Bolshevik strategy in the tsar's parliament, ``For so long as workers and their revolutionary organizations are not strong enough to overthrow parliament\ldots it can be used as a platform for revolutionary socialist propaganda.\ldots\ For this, of course, the center of gravity of the party should be outside parliament.''\footnote{Tony Cliff, ``Introduction,'' in A.\ Y.\ Badeyev, \emph{Bolsheviks in the Tsarist Duma} (Chicago: Bookmarks, 1987), 7.}

There was an understandable weariness of opportunism among the German communists after the experience with the SPD. But while the German SPD had made its ministers of parliament the leaders of the party, Lenin argued that revolutionaries must subordinate parliamentary activity to the overall class struggle, not the other way around:

\blockquote{Criticism---the keenest, most ruthless and uncompromising criticism---must be directed, not against parliamentarism or parliamentary activities, but against those leaders who are unable---and still more against those who are \emph{unwilling}---to utilize parliamentary elections and the parliamentary tribune in a revolutionary, Communist manner.\footnote{Lenin, 48.}}

\paragraph{``Get out of the unions!''}

The KPD wasn't content to simply sit out the elections. Many members also wanted to abstain from Germany's swelling labor unions---most often led by the SPD---as well. As some of the German ``lefts'' argued:

\blockquote{New forms of organization must be created.\ldots\ A Workers' Union, based on factory organizations, should be the rallying point for all revolutionary elements. This should unite all workers who follow the slogan, ``Get out of the trade unions.''\footnote{Hallas, 40.}}

Precisely while some in the KPD were telling workers to get ``Out of the Unions!'' many workers, for the first time, were flooding into them. This was a step toward a more organized and class-conscious working class. That these unions were largely controlled by the SPD was incontestable. The SPD, now in government, was also persecuting strikers and communists. The German ``lefts'' argued that in their new ``Workers Union,'' the ``recognition of the class struggle, of the Soviet system and the dictatorship of the proletariat should be sufficient for enrolment.''\footnote{Ibid.} But, Lenin argued, the ultra-lefts had things backwards:

\blockquote{Greater foolishness and greater damage to the revolution than that caused by such ``Left'' revolutionaries cannot be imagined! Why, if we in Russia today, after two and half years of unprecedented victories over the bourgeoisie of the Entente, were we to make ``recognition of the dictatorship'' a condition of trade union membership, we should be committing a folly, we should be damaging our influence over the masses and should be helping the Mensheviks [the Russian reformists].\footnote{Lenin, 38.}}

The fact that there was reformist, even reactionary influence in the trade unions, meant that revolutionaries should work inside them and not cede ground to such forces. Calling for workers to get out of the unions, for immaculate revolutionary unions, was to write off the millions of workers who were open, but not yet won, to the need for revolution. But inside those unions, Lenin argued, the KPD could expose the vacillation and timidity of the SPD leaders:

\blockquote{If you want to\ldots win the sympathy, confidence and support of ``the masses,'' you must not fear difficulties, you must not fear the pin-pricks, chicanery, insults and persecution of the ``leaders'' (who, being opportunists and social chauvinist, are in most cases directly or indirectly connected with the bourgeoisie and the police), but must imperatively work \emph{wherever the masses are to be found}.\footnote{Ibid., 37.}}

\paragraph{``Our theory is not a dogma but a guide to action''}

Many in the KPD who opposed working in reactionary parliaments and trade unions unsurprisingly objected to making any compromises whatsoever. Again, there was an element of healthy distrust of compromising and temporizing after the betrayals of the social democrats. But it meant, short of a made-to-order pristine revolution, abstaining from the class struggle. The Frankfurt Lefts, for example, rejected ``most emphatically all compromises with other parties\ldots all policy of maneuvering and compromise.''\footnote{Quoted in Ibid., 49.}

In fact, Lenin argued, that would mean rejecting Bolshevism as well. In order to survive the vicissitudes of underground work, a revolution (in 1905), reaction, a strike wave, the patriotic frenzy that accompanied the arrival of the First World War, and the revolution that ended it, the Bolshevik Party had to, as Lenin put it, maneuver, temporize, and compromise:

\blockquote{To carry on a war for the overthrow of the international bourgeoisie, a war which is a hundred times more difficult, protracted and complex than the most stubborn of ordinary wars between states, and refuse before hand any change of tack, or any utilization of a conflict of interests (even if temporary) among one's enemies, or any conciliation or compromise with possible allies (even if they are temporary, unstable, vacillating or conditional allies)---is that not ridiculous in the extreme? Is it not like making a difficult ascent of an unexplored and hitherto inaccessible mountain and refusing in advance ever to move in zigzags, ever to retrace one's steps, or ever to abandon a course once selected, and to try others?\footnote{Lenin, 52.}}

It is not enough, Lenin wrote, to know how to advance. Revolutionary parties also have to know ``how to retreat properly.''\footnote{Ibid.} To look at the actual terrain of the mountain, and chart the best possible course to the top as they go along. Giant boulders can't be wished away. They must be dealt with. After the defeat of the 1905 Revolution, Lenin wrote that the Bolsheviks made the most effective retreat, with ``the least loss to their `army,' with its nucleus best preserved, with the least demoralization, and the best condition to resume the work on the broadest scale and in the most correct and energetic manner.''

The Bolsheviks achieved this only because they ruthlessly exposed and expelled the revolutionary phrase-mongers, who refused to understand that one had to retreat, that one had to know how to retreat, and that one had absolutely to learn how to work legally in the most reactionary parliaments, in the most reactionary trade unions, cooperative societies, mutual insurance and similar organizations.\footnote{Ibid., 13.}

Being able to work inside trade unions and parliaments is not the same as abandoning the fight, or \emph{conforming} to the trade union leadership, or becoming enamored by parliamentary posts. There is ``opportunism''---adapting to the movement as it is---and working inside the labor movement to push it leftward. As Lenin put it:

\blockquote{Capitalism would not be capitalism if the ``pure'' proletariat were not surrounded by a large number of exceedingly mixed transitional types.\ldots\ And all this makes it necessity, absolutely necessary, for the vanguard of the proletariat, its class-conscious section, the Communist Party, to resort to maneuvers, arrangements and compromises with the various groups of proletarians, with the various parties of the workers and small proprietors. The whole point lies in knowing how to apply these tactics in such a way to raise---and not lower---the general level of proletarian class consciousness, revolutionary spirit, and ability to fight and conquer.\footnote{Ibid, 13--14.}}

Likewise, Marxists could not always advance, onward and upward. The Brest-Litovsk Peace Treaty, signed by the Russian workers' government with Germany, was an example of a necessary retreat and compromise. Although huge amounts of Russian territory were (reluctantly) signed away, it was necessary to end the war and give the revolutions in the West time to catch up. Lenin was forced to argue with a number of Russian Bolsheviks who wanted to wage a revolutionary war with Germany. Aside from the fact it would be difficult to prove the solidarity of Russian workers with German workers while the fighting continued, Russian workers had no resources to wage such a war. The seemingly more radical option would have proven a disaster.

In a sense, all labor agreements are compromises in which the terms of exploitation, not exploitation itself, are negotiated between labor and capital. Clearly, not every strike can become a revolution. It must at some point end in an agreement that balances towards one side or the other. The question is, does the outcome of the class struggle at any given stage advance the position of the workers' involved, or set it back? If the workers are forced to retreat, is the retreat forced by circumstances, or based on mistakes or bureaucratic misleadership?

\paragraph{The British Labour Party}

Tony Blair's Labour Party in Britain inspires little (aside from perhaps contempt) these days. But the Labour Party in 1920, was something new, and was a conundrum for British revolutionaries. Workers had only just broken from the Liberal Party (who along with the Tories gave Britain two capitalist parties, much as in the U.S.\ today). As Lenin remarked at the Second Congress of the Communist International:

\blockquote{[The] British Labour Party is in a very special position: it is a highly original kind of party, or rather it is not a party in the ordinary sense of the word. It is made up of members of all trade unions, and has a membership of about four million, and allows sufficient freedom to all affiliated political parties. It thus includes a vast number of British workers who follow the lead of the worst bourgeois elements\ldots however, the Labour Party has let the British Socialist Party into its ranks, permitting it to have its own press organs in which members\ldots can declare that the party leaders are social traitors\ldots a party which unites enormous masses of workers\ldots is nevertheless obliged to grant its members complete latitude.\ldots\ In such circumstances it would be a mistake not to join such a party.\footnote{Hallas, 43.}}

In \emph{Left-Wing Communism}, Lenin had left the question of the British communists affiliating with the Labour Party open. However he did criticize comrades who were entirely dismissive of the idea. If there were millions of workers looking to this party, and it was possible to get in there with them, without sacrificing the right to organize and make independent propaganda, why not do so? But at the time, a Scottish communist Gallacher wrote, ``any support given to parliamentarism is simply assisting to put power into the hands of our British Scheidemanns and Noskes.''\footnote{Lenin, 61.}

Lenin observed that this displayed ``the germs of all the mistakes that are being committed by the German `Left' Communists and that were committed by the `Left' Bolsheviks in 1908 and 1918.''\footnote{Ibid.} Having lost their two bosses' parties, the liberal capitalists were suddenly looking to the Labour Party to save them, as it enjoyed the support of most workers. But British communists were convinced that they ``must not dissipate our energy in adding to the strength of the Labour Party; its rise to power is inevitable. We must concentrate on making a Communist movement that will vanquish it.''\footnote{Ibid., 65.}

Again, for Lenin, the ``ultra-lefts'' turned things on their head. The forces with which they could ``vanquish'' the Labour Party were inside the Labour Party. But how to win them? The fact that the Party had not yet proven itself to be treacherous in power, that the liberals were hijacking it, and that most of the working class looked to it, were all reasons to try to affiliate. Not to hide one's own politics, but to help show greater numbers of workers the true colors of the reformist leaders.

Lenin's approach to affiliation was purely tactical. If it was possible to affiliate and retain full freedom of criticism and organization, then affiliation was useful, because it would put British communist workers in touch with the majority of reformist workers.\footnote{The point quickly became moot when the Labour Party denied the newly-formed British Communist Party's application for affiliation.} For British ultra-leftists like Sylvia Pankhurst, on the other hand, such questions as affiliation, or even participation in electoral politics, were rejected \emph{on principle}. She wrote, for example, ``The Communist Party must keep its doctrine pure, and its independence of reformism inviolate, its mission is to lead the way, without stopping or turning, by the direct road to the communist revolution.''\footnote{Lenin, 65.}

Lenin argued that this position simply abandoned the workers to their reformist leaders:

\blockquote{[The] fact that the majority of the workers still follow the lead of the British Kerenskys or Scheidemanns [Russian and German reformist party leaders, respectively] and that they have not had a government composed of these people, which\ldots was required in Russia and Germany to secure the mass passage of the workers to Communism\ldots indicates that the British Communists \emph{should} participate in parliamentary action, that they should from \emph{within} Parliament help the masses of the workers to see the results of a Henderson and Snowden [Labour Party leaders] government in practice.\footnote{Ibid.}}

\paragraph{Imaginary human material}

Far from the liberal (and Stalinist) myths of his being a dictatorial Svengali, Lenin's starting and ending point, for nearly his entire adult political life, was the idea that the ``emancipation of the working class could only be achieved by the working class itself.'' That is why he rejected both the idea that a minority acting on behalf of the workers could achieve socialism, or that socialists should seek popularity by pandering to any backwardness among workers:

\blockquote{We can (and must) begin to build Socialism not with imaginary human material, not with human material invented by us, but with the human material bequeathed to us by capitalism. That is very ``difficult,'' it goes without saying, but no other approach to the task is serious enough to warrant discussion.\footnote{Ibid., 340.}}

Workers in the U.S.\ today still vote for Republicans and Democrats. Some workers have racist, homophobic, and sexist ideas. The labor movement is still struggling to find its legs after years of atrophy. Here and there, on a small scale, workers are just beginning to fight back. As Duncan Hallas once wrote, to state facts such as these ``is sometimes regarded as something of a betrayal, a slander against the working class. And yet it is merely a statement, not only of what exists, but also of what must exist for capitalist class society in its `democratic' form to continue at all.''\footnote{Duncan Hallas, ``Towards a Revolutionary Party,'' \emph{International Socialist Review} 24, July--August, 2002, 7.} The question is what role revolutionaries can play in breaking away, depending on the level of struggle and crises, layers of workers from supporting the system that exploits and oppresses them. This cannot be done by barking condemnations from the sidelines, or abstaining from real struggles because they are not ``radical'' enough, but by fighting back alongside workers for immediate demands, however small, and in the process winning them to socialist politics.

Neither are such facts an excuse to give up on the revolutionary potential of the working class. Too often have socialists given up on the working class only in time to witness an explosion of struggles and revolutions. The question is to figure out what will help increase the chances of success in both the immediate struggle and in the ultimate fight. As Rosa Luxemburg wrote, Lenin and the Bolsheviks were among the first to deal with the ``practical placing of the problem of the realization of socialism.''\footnote{Rosa Luxemburg, ``The Russian Revolution,'' in \emph{The Rosa Luxemburg Reader} (New York: Monthly Review Press, 2004), 310.} And that is why they still have so much to teach us today.
